\section{A2C-Nachbau auf neuen Daten (Boeing)}

\textbf{Protokoll (aus dem Notebook übernommen):} fünf Agenten mit dem Algorithmus
\texttt{A2C} (Policy \texttt{MlpPolicy}), 75.000 Trainingsschritte je Agent, Auswertung mit
\texttt{ai\_trade\_performance} (\texttt{num\_plays=20}). Der Test-Environment übernimmt die
Merkmalsliste (\texttt{use\_variables}) und den Skalierer des Trainings-Environments. Agent und
Buy-and-Hold laufen auf denselben zwanzig Episoden-Seeds, der Vergleich ist also fair.

\begin{table}[H]
  \centering
  \caption{Zeitfenster des Nachbaus.}
  \label{tab:fenster}
  \begin{tabular}{lll}
    \toprule
     & Zeitraum & Datenpunkte\\
    \midrule
    Trainingsfenster (in-sample)   & 06.10.2021 -- 31.05.2024 & ca.\ 660 Handelstage\\
    Testfenster (out-of-sample)    & 01.06.2024 -- 06.10.2026 & ca.\ 580 Handelstage\\
    \bottomrule
  \end{tabular}
\end{table}

\subsection{Ergebnisse je Agent}

End-NAV je Episode (Start $=1{,}0$), Mittel über zwanzig Episoden pro Agent.

\begin{table}[H]
  \centering
  \caption{2026-Nachbau auf Boeing: End-NAV, Mittel über zwanzig Episoden.}
  \label{tab:a2c}
  \begin{tabular}{crrr|rrr}
    \toprule
    \# & Seed & Training Agent & Training B\&H & Test Agent & Test B\&H & Agent besser\\
    \midrule
    1 & 100 & 1,6705 & 0,9647 & \textbf{0,5616} & 0,9475 & 0\,\%\\
    2 & 101 & 2,4073 & 0,9671 & 1,2403 & 0,9408 & 100\,\%\\
    3 & 102 & 4,0534 & 0,9956 & 0,9828 & 0,9412 & 75\,\%\\
    4 & 103 & 7,0580 & 0,9711 & 1,0106 & 0,9493 & 70\,\%\\
    5 & 104 & 5,1547 & 0,9727 & 1,1326 & 0,9485 & 80\,\%\\
    \midrule
    $\varnothing$ & & 4,0688 & 0,9742 & 0,9856 & 0,9455 & 65\,\%\\
    \bottomrule
  \end{tabular}
\end{table}

\subsection{Handelsverhalten}

\begin{table}[H]
  \centering
  \caption{Aktionsverteilung der Agenten (je eine Episode).}
  \label{tab:aktionen}
  \begin{tabular}{crrr|rrr}
    \toprule
    \# & Training short & Training cash & Training long & Test short & Test cash & Test long\\
    \midrule
    1 & 268 & 18  & 348 & 239 & 17  & 299\\
    2 & 295 & 39  & 300 & 304 & 24  & 227\\
    3 & 297 & 24  & 313 & 279 & 29  & 247\\
    4 & 246 & 212 & 176 & 260 & 201 & 94\\
    5 & 296 & 15  & 323 & 309 & 12  & 234\\
    \bottomrule
  \end{tabular}
\end{table}

Alle Agenten handeln überwiegend long; nur Agent~4 hält eine große Cash-Position. Short-Positionen
werden fast immer gehalten, aber selten stark gewichtet.

\subsection{Interpretation}

\paragraph{In-sample ist kein Beleg für Können.} Im Trainingsfenster erreichen die Agenten im
Mittel einen NAV von 4,07 gegenüber 0,97 für Buy-and-Hold und schlagen die Benchmark in
\emph{jeder} Episode. Das ist kein Ergebnis: Hier wurde trainiert und hier wurde gemessen, die
Zahl zeigt nur, dass die Agenten das Fenster auswendig gelernt haben.

\paragraph{Out-of-sample ist das Bild nüchtern.} Im Testfenster liegt der Mittelwert bei
0,9856 gegenüber 0,9455 für Buy-and-Hold, der Agent ist in \textbf{65\,\%} der Episoden
besser -- aber mit sehr großer Streuung: Agent~1 verliert deutlich (0,56), Agent~2 gewinnt
klar (1,24), und zwei von fünf Agenten verlieren in mindestens einem Viertel der Episoden
gegen Buy-and-Hold. Der Vorteil liegt innerhalb der Seed-Streuung und ist
\textbf{nicht signifikant}.

\paragraph{Buy-and-Hold liegt in beiden Fenstern unter 1,0}, Boeing hat also in beiden
Zeiträumen an Wert verloren. Die Frage lautet hier \glqq weniger verlieren\grqq{} statt
\glqq gewinnen\grqq. Insgesamt \textbf{stützt} der Nachbau die Konklusion des Original-Notebooks:
A2C liefert keine verlässliche Outperformance.

\begin{figure}[H]
  \centering
  \includegraphics[width=0.98\textwidth]{update_2026_a2c.png}
  \caption{End-NAV je Episode (Mittel der fünf Agenten), links Trainingsfenster
  (in-sample), rechts Testfenster (out-of-sample).}
\end{figure}
